\notechapter{N-20220721}{Probability, Conditional Probability, and Substitution Tricks in Inequalities}

\section{Random experiments and sample spaces}

This note begins with the language of probability.  A random experiment is a process that can be repeated under the same conditions, whose possible outcomes are known, but whose actual outcome is uncertain.  Examples include tossing a coin, drawing a ball from a box, or choosing a card.

The set of all possible basic outcomes is the sample space, denoted
\[
  \Omega.
\]
A subset \(A\subseteq\Omega\) is an event.  If the sample space is finite and all outcomes are equally likely, then
\[
  P(A)=\frac{|A|}{|\Omega|}.
\]

The note stresses that events are sets.  This is why union, intersection, complement, and inclusion--exclusion immediately become probability tools:
\[
  P(A\cup B)=P(A)+P(B)-P(A\cap B).
\]

\section{Classical model and Venn diagrams}

One example has seven labeled balls, with three properties represented by sets \(A_1,A_2,A_3\).  The note uses a Venn diagram and factorial counts to compute a probability such as
\[
  P=\frac{7!-3\cdot6!+3\cdot5!-4!}{7!}.
\]
The exact numbers matter less than the method.  If conditions overlap, count by inclusion--exclusion.  A Venn diagram is not a proof by itself, but it helps prevent double-counting.

\section{Conditional probability}

The conditional probability of \(B\) given \(A\) is
\[
  P(B\mid A)=\frac{P(A\cap B)}{P(A)}
\]
when \(P(A)>0\).  The note includes examples where one must be careful not to confuse \(P(B\mid A)\) with \(P(B)\).

Independence is the condition
\[
  P(A\cap B)=P(A)P(B),
\]
or equivalently
\[
  P(B\mid A)=P(B).
\]
This means that knowing \(A\) occurred does not change the probability of \(B\).

A common trap appears in nested events.  If \(A\subseteq B\), then
\[
  P(B\mid A)=1,
\]
but
\[
  P(A\mid B)=\frac{P(A)}{P(B)}
\]
may be much smaller.  Conditional probability is not symmetric.

\section{A transition to substitution tricks}

The note then changes direction and discusses substitutions used in inequalities.  A recurring constraint is
\[
  abc=1.
\]
A standard substitution is
\[
  a=\frac xy,\qquad b=\frac yz,\qquad c=\frac zx.
\]
Another useful substitution is
\[
  x=a+\frac1a,\qquad y=b+\frac1b,\qquad z=c+\frac1c.
\]
Such substitutions are not just algebraic decoration.  They encode hidden positivity and symmetry.  The note points out that many complicated-looking inequalities become simpler after this change of variables.

\section{A classical symmetric equation}

The pages discuss equations such as
\[
  xyz=x+y+z+2.
\]
This can be rewritten as
\[
  xyz+1=(x+1)(y+1)(z+1)-xy-yz-zx-x-y-z.
\]
A better-known substitution is
\[
  x=a+\frac1a,\quad y=b+\frac1b,\quad z=c+\frac1c,
\]
or trigonometric substitutions when the variables resemble
\[
  2\cos A,\quad 2\cos B,\quad 2\cos C.
\]
The note is circling around the idea that symmetric algebraic equations often become geometry or trigonometry after the right substitution.

\section{Inequality examples}

One example asks to prove
\[
  \frac{a}{b+c+a}+\frac{b}{c+a+b}+\frac{c}{a+b+c}\ge\frac32
\]
in a normalized form.  The exact expression on the page is a contest inequality, and the solution uses standard tools such as Cauchy, AM--GM, or substitution.

The point is that conditions like \(abc=1\), \(x+y+z=\text{constant}\), or \(x^2+y^2+z^2+2xyz=1\) should not be treated as random constraints.  They are invitations to choose a special parametrization.

\section{Jensen, trigonometry, and contest style}

Later examples involve trigonometric inequalities in a triangle.  A triangle condition \(A+B+C=\pi\) naturally suggests identities such as
\[
  \cos^2A+\cos^2B+\cos^2C+2\cos A\cos B\cos C=1.
\]
This identity is one of the reasons that substitutions by cosines appear in olympiad inequalities.

\section{The hidden structure}

This note looks like probability followed by inequalities, but both topics are about choosing the right sample space.  In probability, the sample space is literal: events are subsets of \(\Omega\).  In inequalities, the ``sample space'' is the domain allowed by the constraint.  A good substitution changes that domain into something more natural, just as choosing a good probability model turns a counting problem into a set problem.

\section{Conditional probability as changing measure}

Conditional probability is not just a formula.  It changes the universe from \(\Omega\) to the event \(B\):
\[
  \mathbb P(A\mid B)=\frac{\mathbb P(A\cap B)}{\mathbb P(B)}.
\]
The denominator renormalizes the measure so that \(B\) has probability \(1\).

This is why Venn diagrams help: conditioning restricts attention to one region and rescales its area.

\section{Bayes formula and inverse problems}

Bayes' formula
\[
  \mathbb P(A_i\mid B)=\frac{\mathbb P(B\mid A_i)\mathbb P(A_i)}{\sum_j \mathbb P(B\mid A_j)\mathbb P(A_j)}
\]
turns likelihoods into posterior probabilities.  It is a finite version of inference: update prior weights after observing evidence.

The formula is especially useful when \(B\) is easier to compute conditional on each case \(A_i\) than directly.

\section{Inequality substitutions as coordinate choices}

The substitution tricks in the inequality part should be read as coordinate changes.  Symmetric expressions often simplify under
\[
  u=x+y,\qquad v=xy,
\]
or under trigonometric substitutions that encode constraints.  The good substitution is the one that turns the constraint into a simple domain.

\section{Changing measures and changing coordinates}

Probability and inequality substitutions share a habit: change the viewpoint before computing.  Conditioning changes the measure; Bayes reverses conditional information; substitutions choose coordinates adapted to the constraint.
